\chapter{Sprint 7: Evaluation and Optimization}
\chaptertoc

\section{Introduction}
Sprint 7 acted as the validation pass for the full stack. The work checked whether perception, API, frontend, and alerting held together under tests, retries, and throughput pressure.

\section{Sprint Backlog}
Sprint 7 backlog included:
\begin{itemize}
	\item consolidate backend and frontend automated test coverage,
	\item validate feed lifecycle, stream continuity, and API contracts,
	\item optimize runtime throughput and responsiveness,
	\item harden failure handling (network drops, stale data, worker exits),
	\item document residual limitations for final delivery planning.
\end{itemize}

\section{Design}

\subsection{Test Plan}
Validation followed a layered strategy:
\begin{itemize}
	\item \textbf{Backend unit/integration tests:} queue analytics logic, API endpoints (including transport capability and WebRTC signaling), runtime streaming behavior, webhook payload handling.
	\item \textbf{Frontend integration tests:} dashboard setup flow, websocket event handling, WebRTC lifecycle/fallback behavior, and API error handling.
	\item \textbf{Scenario tests:} n8n send/acknowledge flows, unauthorized branch, and multi-feed alert scenarios.
\end{itemize}

\subsection{Performance Metrics}
The optimization targets were tied directly to operational requirements:
\begin{itemize}
	\item processing responsiveness (loop cadence and processed-frame throughput),
	\item stream continuity under feed state transitions,
	\item webhook impact on main-loop latency,
	\item dashboard freshness under websocket degradation.
\end{itemize}

\section{Implementation}
\subsection{Automated Test Coverage}
The test suite covers API transport-capability and WebRTC-offer cases, runtime stream behavior, webhook delivery handling, dashboard integration, live-state synchronization, API error contracts, and transport fallback behavior.
These tests verify API envelopes, event contracts, WebRTC signaling and fallback paths, stream handoff, and setup-workflow behavior.

\subsection{Runtime Optimization Actions}
Key optimizations implemented across the runtime loop and configuration layers include:
\begin{itemize}
	\item configurable detector image size and frame stride,
	\item explicit inference device control and diagnostics,
	\item asynchronous webhook dispatch to keep the processing loop responsive,
	\item real-time playback pacing for file sources.
\end{itemize}

\subsection{Resilience Improvements}
Runtime resilience was improved at multiple levels:
\begin{itemize}
	\item feed worker supervision and exit recovery in \texttt{FeedRegistry},
	\item capture-to-worker frame handoff continuity in \texttt{FeedFrameStreamManager},
	\item websocket reconnect/fallback logic on the frontend side.
\end{itemize}

\subsection{Observed Outcome}
By the end of the sprint, the system had reached a stable integration point where perception, analytics, API runtime, dashboard, and n8n automation operated as one pipeline. The remaining work is mostly benchmarking and deployment hardening rather than missing core functionality.

\section{Conclusion}
Sprint 7 closed the loop with structured testing and targeted performance improvements. The resulting platform is solid enough for demonstration and academic evaluation, with a clear path to production hardening.
